\section{Расчётная задача 5}

Найти собственные числа и координаты собственных векторов линейного оператора, имеющего в некотором базисе матрицу
\[
	A :=
	\begin{pmatrix}
		4 & -3 & -3 \\
		1 & 2 & 1 \\
		1 & 1 & 2
	\end{pmatrix}
	.
\]
Если существует базис пространства, состоящий из собственных векторов данного оператора, требуется также состоавить матрицу оператора в этом базисе.

\begin{proofm}
	Пусть~$V$ --- векторное пространство над полем~$\kk$, $\ee := (e_1, e_2, e_3) \subset V$ --- базис, в котором оператор~$\phi \colon V \to V$ имеет заданную матрицу~$A$.

	Составим характеристический многочлен~$\chi_\phi(t) \in \kk[t]$ оператора~$\phi$:
	\[
		\chi_\phi(t) :=
		\det(t \id - \phi)
		.
	\]
	(Можно было бы вычислить $\det(\phi - t\id)$ вместо~$\det(t\id - \phi)$.)
	Собственные числа оператора~$\phi$ (его \tdef{спектр}) --- это множество нулей~$\chi_\phi(t)$. Этот многочлен инвариантен относительно замены базиса, поэтому его можно вычислить в фиксированной паре базисов~$(\ee, \ee)$, в которой~$\phi$ имеет матрицу~$A$, а~$\id$ --- единичную матрицу:
	\begin{multline}
		\chi_\phi(t) =
		\chi_A(t) =
		\det(tE - A) =
		\det
		\left[ 
			t
			\begin{pmatrix}
				1 & 0 & 0 \\
				0 & 1 & 0 \\
				0 & 0 & 1
			\end{pmatrix}
			-
			\begin{pmatrix}
				4 & -3 & -3 \\
				1 & 2 & 1 \\
				1 & 1 & 2
			\end{pmatrix}
		\right]
		\nline{=}
		\det
		\begin{pmatrix}
			t-4 & 3 & 3 \\
			-1 & t-2 & -1 \\
			-1 & -1 & t-2
		\end{pmatrix}
		=
		(-1)^3
		\det
		\begin{pmatrix}
			t-4 & 3 & 3 \\
			-1 & t-2 & -1 \\
			-1 & -1 & t-2
		\end{pmatrix}
		.
		\label{eq:num5:chi}
	\end{multline}
	Раскрыть опеределитель можно, например, по правилу Саррюса (при этом во многих случаях бывает выгоднее не раскрывать скобки сразу):
	\begin{align*}
		\raisebox{0pt}[2em][0pt]{$
			\displaystyle
			\det
			\begin{pmatrix}
				t-4 & 3 & 3 \\
				-1 & t-2 & -1 \\
				-1 & -1 & t-2
			\end{pmatrix}
		$}
		&=
		(t-4)(t-2)(t-2) +
		{} \\ {} &\quad +
		3 \cdot (-1) \cdot (-1) +
		{} \\ {} &\quad +
		3 \cdot (-1) \cdot (-1) -
		{} \\ {} &\quad -
		(-1) \cdot (t - 2) \cdot 3 -
		{} \\ {} &\quad -
		(-1) \cdot 3 \cdot (t-2) -
		{} \\ {} &\quad -
		(-1) \cdot (-1) \cdot (t-4)
		\nlinea{=}
		(t-4)(t-2)^2 +
		6 +
		6(t-2) - (t-4)
		\nlinea{=}
		(t-4)(t-2)^2 +
		5t - 2
		.
	\end{align*}
	Теперь найдём его корни, то есть решим уравнение~$\chi_\phi(t) = 0$. Множитель~$(-1)^3 = -1$ перед~определителем в~\eqref{eq:num5:chi} не влияет на решение этого уравнения.
	\begin{align*}
		(t^2  - 4t + 4)(t - 4) + 5t - 2
		&= 0
		\\
		t^3 - 4t^2 + 4t - 4t^2 + 16t - 16 + 5t - 2
		&= 0
		\\
		t^3 - 8t^2 + 25t - 18
		&= 0
		.
	\end{align*}
	Заметим, что~$t = 1$ является корнем этого уравнения. Поэтому поделим~$\chi_\phi(t)$ на моном~$(t-1)$. Получим $t^2 - 7t + 18$. Тогда уравнение примет вид
	\[
		(t-1)(t^2 - 7t + 18) = 0.
	\]
	Находим корни квадратичного полинома:
	\[
		t
		=
		\frac{7 \pm \sqrt{49 - 4 \cdot 18}}{2}
		=
		\frac{7 \pm \sqrt{49 - 72}}{2}
		=
		\frac{7 \pm i\sqrt{23}}{2}
		.
	\]
	Видно, что он неприводим над~$\RR$ (у него нет вещественных корней). Поэтому у~$\phi$ как оператора в пространстве~$V$ \emph{над~$\RR$} есть одно инвариантное подпространство размерности~$\deg(t-1) = 1$, соответствующее собственному числу~$\lambda_0 = 1$, и одно инвариантное подпространство размерности~$\deg(t^2 - 7t + 18) = 2$, которое соответствует собственным числам
	\[
		\lambda_\nu =
		\frac{
			7 + (-1)^\nu i\sqrt{23}
		}{2}
		,
		\qquad
		\nu \in \set*{1, 2}
		.
		\proofmark
	\]
\end{proofm}
